\section*{NeurIPS Paper Checklist}

\begin{enumerate}

\item {\bf Claims}
    \item[] Question: Do the main claims made in the abstract and introduction accurately reflect the paper's contributions and scope?
    \item[] Answer: \answerYes{}.
    \item[] Justification: The abstract and introduction frame the contribution as a dataset and benchmark protocol for grounded Indian legal QA, with conservative claims tied to the reported retrieval and generation results.

\item {\bf Limitations}
    \item[] Question: Does the paper discuss the limitations of the work performed by the authors?
    \item[] Answer: \answerYes{}.
    \item[] Justification: The paper includes a limitations and responsible use section covering dataset coverage, privacy, bias, leakage, metric limitations, and legal-advice misuse.

\item {\bf Theory assumptions and proofs}
    \item[] Question: For each theoretical result, does the paper provide the full set of assumptions and a complete (and correct) proof?
    \item[] Answer: \answerNA{}.
    \item[] Justification: The paper does not present theoretical results or proofs; it introduces a dataset, benchmark protocol, and empirical baselines.

\item {\bf Experimental result reproducibility}
    \item[] Question: Does the paper fully disclose all the information needed to reproduce the main experimental results of the paper to the extent that it affects the main claims and/or conclusions of the paper (regardless of whether the code and data are provided or not)?
    \item[] Answer: \answerYes{}.
    \item[] Justification: The paper describes the dataset, retrieval candidate sizes, model, hardware, metrics, and released benchmark scripts; the supplemental material should include the exact commands used for each run.

\item {\bf Open access to data and code}
    \item[] Question: Does the paper provide open access to the data and code, with sufficient instructions to faithfully reproduce the main experimental results, as described in supplemental material?
    \item[] Answer: \answerYes{}.
    \item[] Justification: The dataset is hosted on Hugging Face under Apache-2.0 and benchmark code is prepared for release with local CUDA and Kaggle instructions. If using double-blind review, anonymize code access or use the E\&D single-blind option.

\item {\bf Experimental setting/details}
    \item[] Question: Does the paper specify all the training and test details (e.g., data splits, hyperparameters, how they were chosen, type of optimizer) necessary to understand the results?
    \item[] Answer: \answerYes{}.
    \item[] Justification: The experiments are evaluation-only baselines. The paper specifies the public train split sample, candidate corpus sizes, TF-IDF retrieval, FLAN-T5-base generation, deterministic decoding, and metrics. No model training optimizer is used.

\item {\bf Experiment statistical significance}
    \item[] Question: Does the paper report error bars suitably and correctly defined or other appropriate information about the statistical significance of the experiments?
    \item[] Answer: \answerNo{}.
    \item[] Justification: The current draft reports single-run baseline results. Before final submission, bootstrap confidence intervals over the 500-example evaluation set should be added if time permits.

\item {\bf Experiments compute resources}
    \item[] Question: For each experiment, does the paper provide sufficient information on the computer resources (type of compute workers, memory, time of execution) needed to reproduce the experiments?
    \item[] Answer: \answerYes{}.
    \item[] Justification: The paper reports an NVIDIA RTX 3050 6GB Laptop GPU and wall-clock runtimes for the main experiments.

\item {\bf Code of ethics}
    \item[] Question: Does the research conducted in the paper conform, in every respect, with the NeurIPS Code of Ethics \url{https://neurips.cc/public/EthicsGuidelines}?
    \item[] Answer: \answerYes{}.
    \item[] Justification: The work is a dataset and evaluation study using public legal materials, with responsible-use limitations and no deployment of automated legal decision-making systems.

\item {\bf Broader impacts}
    \item[] Question: Does the paper discuss both potential positive societal impacts and negative societal impacts of the work performed?
    \item[] Answer: \answerYes{}.
    \item[] Justification: The paper discusses improved access to legal information and reproducible Indian legal NLP research, along with risks such as hallucinated advice, privacy leakage, bias amplification, and overreliance by non-experts.

\item {\bf Safeguards}
    \item[] Question: Does the paper describe safeguards that have been put in place for responsible release of data or models that have a high risk for misuse (e.g., pre-trained language models, image generators, or scraped datasets)?
    \item[] Answer: \answerYes{}.
    \item[] Justification: The paper describes documentation, responsible-use guidance, Croissant metadata, dataset-card warnings, and human-in-the-loop requirements for downstream legal applications. No new model weights are released.

\item {\bf Licenses for existing assets}
    \item[] Question: Are the creators or original owners of assets (e.g., code, data, models), used in the paper, properly credited and are the license and terms of use explicitly mentioned and properly respected?
    \item[] Answer: \answerYes{}.
    \item[] Justification: The dataset page reports Apache-2.0 licensing and credits the dataset creator. FLAN-T5 and other software dependencies are cited. Exact source provenance and terms for upstream legal sources must be completed before submission.

\item {\bf New assets}
    \item[] Question: Are new assets introduced in the paper well documented and is the documentation provided alongside the assets?
    \item[] Answer: \answerYes{}.
    \item[] Justification: The dataset is documented through a dataset card, Croissant metadata, benchmark code, metrics, and responsible-use notes. The paper also lists remaining provenance fields that should be completed.

\item {\bf Crowdsourcing and research with human subjects}
    \item[] Question: For crowdsourcing experiments and research with human subjects, does the paper include the full text of instructions given to participants and screenshots, if applicable, as well as details about compensation (if any)?
    \item[] Answer: \answerNA{}.
    \item[] Justification: The current benchmark does not include a new crowdsourcing or human-subject experiment. If Argilla reviewers were paid annotators, their instructions and compensation details should be added to the appendix.

\item {\bf Institutional review board (IRB) approvals or equivalent for research with human subjects}
    \item[] Question: Does the paper describe potential risks incurred by study participants, whether such risks were disclosed to the subjects, and whether Institutional Review Board (IRB) approvals or equivalent were obtained?
    \item[] Answer: \answerNA{}.
    \item[] Justification: The current benchmark does not involve newly recruited human subjects. If the dataset construction involved human annotators beyond the authors, document whether review was required and how annotator privacy was protected.

\item {\bf Declaration of LLM usage}
    \item[] Question: Does the paper describe the usage of LLMs if it is an important, original, or non-standard component of the core methods in this research?
    \item[] Answer: \answerYes{}.
    \item[] Justification: The paper discloses FLAN-T5-base as an evaluation baseline and notes that model-assisted question-answer generation, if used in dataset construction, must be documented with prompts, model names, and human review procedure.

\end{enumerate}
